% =============================================================================
% APPENDIX D: LITERATURE REVIEW – CLIMATE RISK FOR DATA CENTERS
% Last updated: 2026-03-10
% =============================================================================

\section{Literature Review}
\label{app:literature}

Results are organised along three guiding questions:
(Q1) physical climate risks for data centers and digital infrastructure;
(Q2) insurance, insurability and climate risk for critical infrastructure;
(Q3) quantitative frameworks to integrate climate risk into infrastructure
valuation.

\subsection{Search Strategy and Screening Criteria}

The search process combined a broad semantic query on climate risk for data
centers and AI-related digital infrastructure with structured screening based
on abstracts and, where available, full texts. Papers and reports were
screened in when they met the following criteria:

\begin{itemize}
  \item \textbf{Infrastructure focus:} data centers, digital infrastructure,
        AI infrastructure, or closely related critical infrastructure
        (energy systems, telecommunications, industrial facilities, utilities).
  \item \textbf{Climate risk analysis:} quantitative modelling or analysis of
        physical climate risks (heat, flooding, storms, sea level rise, drought)
        and/or their economic or operational impacts.
  \item \textbf{Publication quality:} peer-reviewed academic papers or
        high-quality institutional reports from central banks, regulators,
        international organisations, or major research institutes.
  \item \textbf{Financial risk assessment:} explicit analysis of
        climate-related insurance, risk transfer, or methods for integrating
        climate risk into investment and financial assessment.
  \item \textbf{Infrastructure scale:} large-scale critical infrastructure
        rather than purely residential or small commercial buildings.
\end{itemize}

For each included study, we extracted information on infrastructure scope,
hazards, risk modelling approach, financial impacts, insurance treatment,
valuation framework, main findings, and acknowledged limitations.

\vspace{0.3cm}
\noindent The remainder of this appendix summarises the evidence base by
guiding question.

% =============================================
% D.1 Q1 – PHYSICAL CLIMATE RISKS
% =============================================

\subsection{Q1 – Physical Climate Risks for Data Centers and Digital Infrastructure}
\label{app:literature_q1}

The peer-reviewed literature directly quantifying physical climate risks for
data centers is extremely limited: only three studies explicitly address data
center siting and operations under climate change. A broader body of work on
analogous infrastructure — power generation, telecommunications, and
electricity grids — provides transferable methodologies. Flooding (riverine,
pluvial and coastal) emerges as the most consistently analysed hazard, with
heat extremes and storms also prominent. Most studies rely on RCP~4.5 and
RCP~8.5 (or SSP-equivalent) scenarios, with horizons extending from the 2030s
to 2100. Limitations repeatedly noted include coarse hazard-map resolution,
generic vulnerability functions not calibrated to local conditions, and
inadequate treatment of indirect impacts, cascading effects, and compound
hazards.

\subsubsection{Direct Data Center Studies}

\begin{table}[H]
\centering
\small
\caption{Direct data center climate risk studies}
\label{tab:app_lit_dc}
\begin{tabular}{@{}p{3.5cm}p{3.2cm}p{7.0cm}@{}}
\toprule
\textbf{Study} & \textbf{Scope} & \textbf{Key contributions and limitations} \\
\midrule
Pasupuleti (2022) &
Data center siting and capacity decisions &
Integrates HEC\textendash HMS hydrologic simulations into data center
location and capacity planning. Uses rainfall–runoff scenarios to estimate
flood probabilities and downtime risk, combined with CloudSim models to
quantify impacts on service availability and financial loss. Demonstrates a
rare linkage from hydrological hazards to operational disruptions; does not
yet include multi-hazard or insurance effects. \\[0.1cm]
Pasupuleti (2024) &
GIS-based siting under flood and climate risk &
Develops a climate- and flood-aware GIS framework for data center site
selection, incorporating riverine and pluvial flooding, sea-level rise, and
extreme precipitation. High-resolution DEMs and multi-criteria decision
analysis produce spatially explicit risk surfaces. Reports a 63\% reduction in
expected annual damage for hazard-aware siting relative to purely economic
criteria; vulnerability remains represented by generic damage functions. \\[0.1cm]
Pasupuleti (2025) &
GIS–AI framework for resilient siting and operation &
Combines climate datasets and machine learning models to build
risk-informed suitability surfaces and predict outage probability and cooling
energy demand. Shows that modest cost premiums for resilient locations can
be offset by reduced disruption and lower energy expenditure over asset
lifetimes. Focuses on siting and operation but does not quantify portfolio-level
financial tail risks. \\
\bottomrule
\end{tabular}
\end{table}

\subsubsection{Analogous Infrastructure with Transferable Methods}

\begin{table}[H]
\centering
\small
\caption{Representative climate risk studies for analogous infrastructure}
\label{tab:app_lit_analogous}
\begin{tabular}{@{}p{3.5cm}p{3.2cm}p{7.0cm}@{}}
\toprule
\textbf{Study} & \textbf{Infrastructure / hazards} & \textbf{Methodological insights} \\
\midrule
Dawson et al. (2018) &
Energy, water, transport, ICT (UK) &
Provides a systems framework for national climate risk assessment across
infrastructure sectors, including ICT. Identifies flooding as the greatest risk
to all sectors and projects a doubling of users reliant on electricity
infrastructure at flood risk under 2~\textdegree C, tripling under 4~\textdegree C.
Highlights interdependence between energy and ICT; notes limited asset-level
performance data. \\[0.1cm]
Forzieri et al. (2016) &
Critical infrastructure in Europe; heat, drought, coastal floods &
Projects that damages to European critical infrastructures could triple by the
2020s, increase six-fold by mid-century, and exceed 10 times present values
by end of century. Identifies heatwaves, droughts, and coastal floods as
showing the most dramatic rise. Uses pan-European hazard maps and generic
vulnerability functions, with limited sector-specific calibration. \\[0.1cm]
Oughton et al. (2023) &
Global mobile telecommunications (7.6M assets);
tropical cyclones, coastal flooding &
Performs a global vulnerability assessment of mobile telecom infrastructure
using crowdsourced asset data and climate projections under RCP~8.5.
Estimates that a 0.01\% annual probability tropical cyclone event could affect
2.26~million cells (USD~1.01~billion damage, 44\% increase), and coastal
flooding could affect 109.9~thousand cells (USD~2.69~billion damage, 78\%
increase). Demonstrates a scalable methodology directly transferable to
digital infrastructure inventories. \\[0.1cm]
Luo et al. (2023) &
Thermal and hydro power plants; multi-hazard &
Develops a multi-hazard risk framework for power generation projects across
Eastern Europe, the Mediterranean, and Central Asia. Uses publicly available
climate data and technology-specific vulnerability factors to estimate
generation losses under RCP~4.5 and RCP~8.5. Highlights rising water
temperatures as a key overlooked risk; does not explicitly monetise all
operational impacts. \\
\bottomrule
\end{tabular}
\end{table}

\subsubsection{Cross-cutting Limitations}

Across Q1 studies, several limitations are recurrent: coarse resolution of
global hazard maps compared to facility-level siting decisions; reliance on
generic damage functions not calibrated for data center equipment, cooling
systems, or power infrastructure; limited treatment of indirect and
cascading impacts (e.g.\ supply chain disruptions, grid failures); and
under-representation of electric and telecommunications networks relative
to other sectors. Data centers are rarely treated explicitly despite their
growing systemic importance.

% =============================================
% D.2 Q2 – INSURANCE AND INSURABILITY
% =============================================

\subsection{Q2 – Insurance, Insurability and Climate Risk for Critical Infrastructure}
\label{app:literature_q2}

The literature on insurance and insurability in the context of infrastructure
climate risk is sparse, and almost entirely absent for data centers and
digital infrastructure. Of the 25 sources screened for insurance content, only
one study (Kerkhofs et al., 2025) provides a substantive quantitative analysis
of insurance mechanisms and their effect on portfolio tail risks. Most
infrastructure climate risk studies do not address premiums, coverage terms,
market withdrawals, or shifts of risk from private insurers to public budgets.
No work was found that documents protection gaps, premium trajectories, or
coverage restrictions specifically for data centers or comparable high-value
digital assets.

\subsubsection{Available Evidence on Insurance Mechanisms}

\begin{table}[H]
\centering
\small
\caption{Insurance-related findings in reviewed studies}
\label{tab:app_lit_insurance}
\begin{tabular}{@{}p{3.5cm}p{3.2cm}p{7.0cm}@{}}
\toprule
\textbf{Study} & \textbf{Context} & \textbf{Insurance-related contributions} \\
\midrule
Kerkhofs et al. (2025) &
Power firms (India, Thailand) &
Develops an asset-level framework for physical risk assessment and
translation to portfolio-level impacts, including explicit treatment of
insurance. Uses Monte Carlo simulation with spatial correlations and
Value-at-Risk metrics to quantify tail risks. Shows that insurance has a
limited effect on average portfolio losses but materially affects tail risk
distribution. Explores scenarios in which insurers withdraw from high-risk
markets, increasing residual risk for asset owners. \\[0.1cm]
Karagiannakis et al. (2024) &
Power grid fragility &
Discusses use of fragility curves by insurers for risk assessment and
life-cycle costing of grid assets. Does not provide empirical analysis of
premium trends, coverage restrictions, or protection gaps. \\[0.1cm]
Other reviewed studies &
Various infrastructure sectors &
Do not include explicit insurance analysis. Insurance is occasionally
mentioned qualitatively but without quantitative evidence on premiums,
coverage, or protection gaps for critical infrastructure. \\
\bottomrule
\end{tabular}
\end{table}

\subsubsection{Identified Gaps for Data Centers}

For data centers and digital infrastructure, the review did not identify any
peer-reviewed work that:

\begin{itemize}
  \item Quantifies protection gaps (insured vs.\ uninsured losses) at asset or
        portfolio level.
  \item Analyses premium dynamics (trend, volatility, deductibles) for
        climate-exposed facilities.
  \item Documents coverage withdrawals, non-renewals, or market capacity
        constraints for large digital campuses.
  \item Examines how climate-driven insurance conditions affect
        bankability, cost of capital, or asset valuation for data centers.
\end{itemize}

This absence is a critical limitation for any climate-economic framework that
aims to integrate insurance dynamics into data center infrastructure
valuation.

% =============================================
% D.3 Q3 – VALUATION FRAMEWORKS
% =============================================

\subsection{Q3 – Quantitative Frameworks to Integrate Climate Risk into Infrastructure Valuation}
\label{app:literature_q3}

A growing literature proposes quantitative frameworks for integrating climate
risk into infrastructure investment and valuation, but none are tailored to
data centers. Existing approaches include probability-of-default and DSCR
metrics for energy assets, real options analysis for flood protection,
discounted cash flow models with climate adjustments, Expected Annual Damage
coupled with stress testing, and storyline-based frameworks linking climate
scenarios to asset damages. Common elements are scenario-based hazard
projections, hazard-to-impact translation via vulnerability functions, and
financial metrics such as NPV, EAD, or VaR. Key methodological gaps for
application to data centers include the absence of data center-specific
vulnerability functions, incomplete representation of indirect and cascading
impacts, and limited treatment of compound multi-hazard risks.

\subsubsection{Representative Valuation Frameworks}

\begin{table}[H]
\centering
\small
\caption{Representative quantitative valuation frameworks}
\label{tab:app_lit_valuation}
\begin{tabular}{@{}p{3.5cm}p{3.2cm}p{7.0cm}@{}}
\toprule
\textbf{Study} & \textbf{Framework type} & \textbf{Relevance and limitations} \\
\midrule
In et al. (2020, 2021, 2022) &
Probability of default / DSCR for energy infrastructure &
Develop a framework for pricing climate-related risks of energy investments,
using DSCR and probability of default rather than NPV alone. Combine
asset-specific physical and transition risks with financing structure to
estimate timing and size of losses. Demonstrate that renewable projects are
more resilient than fossil assets. Applied only to energy; does not consider
digital infrastructure or supply-chain effects. \\[0.1cm]
Prime et al. (2018) &
Real options for flood protection &
Applies real options analysis to flood defense investment around electricity
infrastructure, using UKCP09 sea-level projections to 2100. Values flexibility
in investment timing under climate uncertainty and shows which sites benefit
from early vs.\ delayed adaptation. Could be adapted for data center flood
defenses; currently focused on coastal electricity assets. \\[0.1cm]
Neumann et al. (2015) &
DCF-based integrated assessment across US infrastructure &
Combines sector models and discounted cash flow (3\% rate) to estimate
climate impacts and adaptation benefits for roads, bridges, coastal
development, and drainage. Finds that mitigation policies reduce cumulative
impacts by 25–35\%, with avoided costs exceeding \$460~billion through 2100.
Electric and telecommunications networks are explicitly excluded. \\[0.1cm]
Pasha et al. (2025) &
Flood risk stress testing (EAD) for solar power &
Implements a climate risk stress test based on the
hazard–vulnerability–exposure approach, using global flood hazard data and
depth–damage functions to derive Expected Annual Damage for a solar plant.
Strength is transparency and replicability; limitations include coarse hazard
resolution, generic damage functions, and omission of indirect financial
impacts. \\[0.1cm]
Koks et al. (2022) &
Storyline framework for coastal infrastructure &
Combines sea-level rise projections, inundation modelling, and infrastructure
damage assessment into a storyline framework. Shows that coordinated
adaptation can reduce coastal flood impacts by more than 70\%. Limited to two
adaptation strategies and coastal flooding, but illustrates how physical
scenarios can be linked to discrete adaptation choices. \\
\bottomrule
\end{tabular}
\end{table}

\subsubsection{Implications for Data Center Valuation}

Taken together, these frameworks demonstrate that climate risk can be
integrated into infrastructure valuation through adjusted DCF, probability-of-default metrics, EAD-based stress testing, and real options. However, no
framework in the reviewed literature has been designed for, or validated on,
data center infrastructure. In particular, existing work does not capture:

\begin{itemize}
  \item Cooling-system dependencies and power quality requirements specific
        to high-density compute.
  \item Operational and business-interruption impacts beyond physical asset
        damage, including service continuity requirements.
  \item Insurance and insurability dynamics as explicit financial channels.
  \item Cascading economic impacts of data center outages on downstream
        users and services.
  \item Multi-hazard compound scenarios (e.g.\ extreme heat, drought, and
        grid stress occurring jointly).
\end{itemize}

